\section{Introduction}

Multilingual pre-trained language models (PLMs) have significantly advanced natural language processing (NLP) by enabling cross-lingual transfer across many languages~\cite{wang2023pretrained,xue2021mt5}. Despite these advances, low-resource languages remain severely underrepresented in training corpora, leading to suboptimal performance on downstream tasks. This underrepresentation is particularly acute for languages that use non-Latin scripts, where standard tokenization strategies allocate disproportionately few subword units, resulting in high out-of-vocabulary (OOV) rates and excessive token fragmentation that degrades performance on token-sensitive tasks such as named entity recognition, sequence labeling, and reading comprehension~\cite{chau2021specializing,ahia2023all}.

A key technical bottleneck is subword tokenization. Standard methods such as Byte Pair Encoding (BPE) and SentencePiece~\cite{kudo2018sentencepiece} are highly effective for high-resource, Latin-script languages but frequently over-segment text in morphologically rich or orthographically distinct languages, producing longer sequences, degraded semantic representations, and elevated OOV rates. Tokens not covered by the model's vocabulary are typically mapped to a generic [UNK] symbol, discarding all lexical information~\cite{liu2024ofa,xue2022byt5}.

These challenges are especially pronounced for Ge'ez-script languages such as Amharic and Tigrinya. The Ge'ez script is a morphosyllabic abugida in which each symbol represents a consonant vowel pair, which gives it a large character inventory~\cite{getatchew1996ethiopic}. Standard multilingual vocabularies built predominantly on Latin-script corpora inevitably produce severe fragmentation for these scripts. Compounding this problem, Amharic and Tigrinya have limited annotated digitized text, orthographic variation across dialects, and inconsistent Unicode encoding in existing datasets, all of which further hamper model adaptation.

Prior work has explored vocabulary expansion for low-resource languages~\cite{wang2019improving}, continued pretraining on parallel corpora~\cite{ebrahimi2021adapt}, adaptation strategies for African languages~\cite{alabi2022adapting}, and large-scale pretraining over 500+ languages~\cite{imanigooghari2023glot500}. However, to our knowledge none of these approaches has been systematically applied to Ge'ez-script languages, and no prior work has combined targeted vocabulary augmentation, language-specific tokenizer construction, mean-based embedding initialization, and two-stage continued pretraining within a unified framework for this language family. Existing broad-coverage models such as Glot500 achieve coverage through scale but do not optimize tokenization for specific script families, leaving a representational gap for Amharic and Tigrinya.

This paper addresses these limitations by proposing \textbf{VEXMLM} (Vocabulary-Extended XLM-R), a systematic framework for adapting XLM-R to Ge'ez-script languages that combines: (1) language-specific SentencePiece tokenization and vocabulary augmentation with \vocabAdded{} Ge'ez-derived subword tokens; (2) mean initialization of new token embeddings to preserve alignment with the existing embedding space; and (3) a two-stage training procedure comprising continued masked language model pretraining followed by task-specific fine-tuning. The adapted model is evaluated on three downstream tasks: named entity recognition (NER) and question answering (QA) in Amharic and Tigrinya, and sentiment analysis (SA) in Amharic.

This work makes three principal contributions:
\begin{itemize}
    \item \textbf{Vocabulary expansion and embedding initialization for Ge'ez script.} We propose a vocabulary augmentation approach, constructing language-specific SentencePiece tokenizers for Amharic and Tigrinya and initializing new token embeddings via the mean of source embeddings, ensuring compatibility with the pretrained XLM-R embedding space.
    \item \textbf{Two-stage training.} We show that continued MLM pretraining on Ge'ez-script monolingual corpora is required for the expanded vocabulary to help: without it, vocabulary expansion lowers Tigrinya NER accuracy on out-of-vocabulary words, and with it the model exceeds the XLM-R baseline.
    \item \textbf{Intrinsic and extrinsic evaluation.} We evaluate VEXMLM with tokenization efficiency metrics (fertility, compression, OOV round-trip) and with downstream performance on NER and QA in Amharic and Tigrinya and sentiment analysis in Amharic, reporting means and standard deviations over five seeds and including the tasks on which VEXMLM does not improve over XLM-R.
\end{itemize}
